% Part: first-order-logic
% Chapter: sequent-calculus
% Section: provability-consistency

\documentclass[../../../include/open-logic-section]{subfiles}

\begin{document}

\iftag{FOL}
      {\olfileid{fol}{seq}{prv}}
      {\olfileid{pl}{seq}{prv}}

\olsection{\usetoken{S}{derivability} en consistentie}

We bewijzen nu een aantal eigenschappen van de relatie
!!{derivability}. Die eigenschappen zijn ook los van elkaar
interessant. Bovendien hebben we ze stuk voor stuk nodig in het bewijs
van de volledigheidsstelling.

\begin{prop}\ollabel{prop:provability-contr}
  Als $\Gamma \Proves !A$ en $\Gamma \cup \{!A\}$ inconsistent is,
  dan is $\Gamma$ inconsistent.
\end{prop}

\begin{proof}
Er zijn eindige verzamelingen $\Gamma_0$ en $\Gamma_1 \subseteq
\Gamma$ waarvoor $\Log{LK}$ de sequenten $\Gamma_0 \Sequent
!A$ en $!A, \Gamma_1 \Sequent \quad$ !!{derive}s. Noem de
$\Log{LK}$-!!{derivation} van $\Gamma_0 \Sequent !A$ nu~$\pi_0$. Noem
de $\Log{LK}$-!!{derivation} van $\Gamma_1, !A \Sequent \quad$
nu~$\pi_1$. Met de snederegel kunnen we dan het volgende !!{derive}:
\begin{prooftree}
\AxiomC{}
\RightLabel{$\pi_0$}
\Deduce$ \Gamma_0 \fCenter !A $
\AxiomC{}
\RightLabel{$\pi_1$}
\Deduce$!A, \Gamma_1 \fCenter $
\RightLabel{\Cut}
\BinaryInf$ \Gamma_0,\Gamma_1 \fCenter $
\end{prooftree}

Omdat $\Gamma_0 \subseteq \Gamma$ en $\Gamma_1 \subseteq \Gamma$, geldt
$\Gamma_0 \cup \Gamma_1 \subseteq \Gamma$. De afleiding hierboven laat
daarom zien dat $\Gamma$ inconsistent is.
\end{proof}

\begin{prop}
\ollabel{prop:prov-incons}
$\Gamma \Proves !A$ precies als $\Gamma \cup \{\lnot !A\}$
inconsistent is.
\end{prop}

\begin{proof}
Neem eerst aan dat $\Gamma \Proves !A$. Er is dan
!!a{derivation}~$\pi_0$ van $\Gamma \Sequent !A$. Voeg daar één stap
met $\LeftR{\lnot}$ aan toe. Zo krijgen we !!a{derivation}
van~$\lnot !A, \Gamma \Sequent \quad$. Dat betekent precies dat
$\Gamma \cup \{\lnot !A\}$ inconsistent is.

Neem nu omgekeerd aan dat $\Gamma \cup \{\lnot !A\}$ inconsistent is.
Dan is er !!a{derivation}~$\pi_1$ van $\lnot !A, \Gamma \Sequent
\quad$. Het volgende is !!a{derivation} van $\Gamma \Sequent !A$:
\begin{prooftree}
  \Axiom$!A \fCenter !A$
  \RightLabel{\RightR{\lnot}}
  \UnaryInf$\fCenter !A, \lnot !A$
  \AxiomC{}
  \RightLabel{$\pi_1$}
  \Deduce$\lnot !A, \Gamma \fCenter$
  \RightLabel{\Cut}
  \BinaryInf$\Gamma \fCenter !A$
\end{prooftree}
\end{proof}

\begin{prob}
Bewijs dat $\Gamma \Proves \lnot !A$ precies als $\Gamma \cup \{!A\}$
inconsistent is.
\end{prob}

\begin{prop}\ollabel{prop:explicit-inc}
  Als $\Gamma \Proves !A$ en $\lnot !A \in \Gamma$, dan is $\Gamma$
  inconsistent.
\end{prop}

\begin{proof}
  Neem aan dat $\Gamma \Proves !A$ en $\lnot !A \in \Gamma$. Dan is er
  !!a{derivation}~$\pi$ van een sequent $\Gamma_0 \Sequent !A$. De
  sequent $\lnot !A, \Gamma_0 \Sequent \quad$ is dan eveneens
  !!{derivable}:
  \begin{prooftree}
    \AxiomC{}
    \RightLabel{$\pi$}
    \Deduce$\Gamma_0 \fCenter !A$
    \Axiom$!A \fCenter !A$
    \RightLabel{\LeftR{\lnot}}
    \UnaryInf$\lnot !A, !A \fCenter$
    \RightLabel{\LeftR{\Exchange}}
    \UnaryInf$!A, \lnot !A \fCenter$
    \RightLabel{\Cut}
    \BinaryInf$\Gamma_0, \lnot !A \fCenter$
  \end{prooftree}
  Omdat $\lnot !A \in \Gamma$ en $\Gamma_0 \subseteq \Gamma$, laat
  deze afleiding zien dat $\Gamma$ inconsistent is.
\end{proof}

\begin{prop}\ollabel{prop:provability-exhaustive}
  Als $\Gamma \cup \{!A\}$ en $\Gamma \cup \{\lnot !A\}$ allebei
  inconsistent zijn, dan is $\Gamma$ inconsistent.
\end{prop}

\begin{proof}
Er zijn dan eindige delen $\Gamma_0 \subseteq \Gamma$ en $\Gamma_1
\subseteq \Gamma$. Verder zijn er $\Log{LK}$-!!{derivation}s $\pi_0$
en $\pi_1$ van respectievelijk $!A, \Gamma_0 \Sequent \quad$ en
$\lnot !A, \Gamma_1 \Sequent \quad$. Met die twee afleidingen kunnen
we het volgende !!{derive}:
\begin{prooftree}
\AxiomC{}
\RightLabel{$\pi_0$}
\Deduce$ !A, \Gamma_0 \fCenter $
\RightLabel{\RightR{\lnot}}
\UnaryInf$ \Gamma_0 \fCenter \lnot !A$
\AxiomC{}
\RightLabel{$\pi_1$}
\Deduce$\lnot !A, \Gamma_1 \fCenter  $
\RightLabel{\Cut}
\BinaryInf$ \Gamma_0, \Gamma_1 \fCenter $
\end{prooftree}
Omdat $\Gamma_0 \subseteq \Gamma$ en $\Gamma_1 \subseteq \Gamma$, geldt
$\Gamma_0 \cup \Gamma_1 \subseteq \Gamma$. De laatste sequent heeft
een lege rechterkant. Dus is $\Gamma$ inconsistent.
\end{proof}
\end{document}
